% Shared main-text method for both active manuscripts.
\section{Preliminaries and Self-Harness Optimization}
\label{sec:method}

An \rlm{} drives a persistent Python REPL to inspect an input, issue model calls,
and assemble an answer~\citep{zhang2025rlm}. The input is initially stored in
\texttt{context}; code can expose selected content through printed output. The
ordinary submission path sets \texttt{answer["content"]} and
\texttt{answer["ready"]}. Model weights, experiment-owned resource caps, and
the environment verifier remain fixed during optimization.

Following \citet{zhang2026selfharness}, we expose ten editable harness surfaces:
five instruction fields, runtime policy, execution-output formatting, REPL
helpers, answer middleware, and a skill library. Each surface corresponds to a
declared field or group of fields. Human-authored capability descriptions and
mechanism-to-surface routes constrain the search. Optimization starts from the
sparse \hzero{}; the hand-authored \hzerostar{} and \texttt{H0*R} serve as
reference harnesses rather than as seeds. Appendix~\ref{app:repro}
lists the frozen artifacts and accounting conventions.

\paragraph{Notation.}
Throughout the paper, $\mathcal{M}$ is the fixed language model, $H_{r-1}$ is
the incumbent entering numbered round $r$ and $H_r$ is its successor,
$D_{\mathrm{in}}$ is the held-in mining split,
$V$ is the repeatedly accessed held-out validation split, and $V_s\subset V$
is its screening subset. Candidate $C_{r,j}$ changes exactly one surface of
$H_{r-1}$; $C_r^\star$ is the candidate selected for full validation.


\subsection{Self-Harness Optimization}
\label{sec:sho}

The same fixed model executes tasks, diagnoses failures, and proposes edits.
Each round has four stages: weakness mining,
proposal of independent candidates, a screening tournament, and full validation
of the tournament winner. Two bounded fallbacks, re-attribution and
refinement, run only when the main path finds no acceptable candidate.

\paragraph{Weakness mining.}
The incumbent $H_{r-1}$ runs on $D_{\mathrm{in}}$. The environment
verifier evaluates final answers, and an attributor uses selected
execution-trace operations to infer failure mechanisms with explicit
uncertainty. Diagnoses use a predefined vocabulary and are clustered by their
recorded failure signature. Cited operations support an interpretation; they
do not independently establish its causal correctness.

\paragraph{Independent candidate proposal.}
The proposer receives compact operation-level evidence, the incumbent surfaces,
runtime capabilities, and bounded history of attempted edits and outcomes,
including mechanisms already tried on other surfaces. It returns up to $k$
candidates $C_{r,1},\dots,C_{r,k}$, each a single-surface edit of $H_{r-1}$ aimed at
one mined pattern. Admission reserves an exact (pattern, surface) pair, so one
pattern may be attacked on several surfaces and one surface may carry edits for
different patterns. Candidates are \emph{not} merged: each is a complete
harness that differs from the incumbent on one surface. For each edit the
proposer must explain what the incumbent does, what observed behavior remains
wrong, and which concrete trigger the edit acts on, phrased as a rule the model
will meet on a real input rather than as softer emphasis. Host materialization
and preflight reject malformed or structurally invalid edits before any task
evaluation.

\paragraph{Screening tournament.}
When more than one candidate is admitted, each one and the incumbent run
independently on a small, fixed screening subset $V_s$ of the validation tasks,
with $v_s<v$ attempts per task. Candidates are ranked by their change in the
promotion metric relative to the incumbent on $V_s$; unscored candidates (for
example, over budget) rank last. Screening only orders candidates: its
accept/reject verdicts never promote a harness. If every screened candidate
falls short of the promotion rule on $V_s$, the loop can re-attribute the
\emph{same} mining runs once, with no new rollouts, proposes a fresh pool, and
screens it. With zero or one admitted candidate the tournament is skipped.

\paragraph{Full validation and promotion.}
Only the top-ranked candidate $C_r^\star$ is fully validated against
$H_{r-1}$ on all held-out validation tasks $V$ with $v$ attempts. Its
screening attempts are reused as part of this evaluation, and an incumbent
that did not change since an earlier round reuses its stored held-out runs.
Let $\bar F_\ell(X)$ be harness $X$'s mean verifier-defined primary quality
(F1 for OOLONG-Pairs, NDCG@10 for OBLIQ) on validation tasks in input-size
stratum $\ell\in\mathcal{L}$ (a context window or a candidate-pool size), and
$\Delta_{r,\ell}=\bar F_\ell(C_r^\star)-\bar F_\ell(H_{r-1})$.
$C_r^\star$ is promoted when
\begin{equation}
\Delta_{r,\ell}\geq-\tau_\ell\ \ \forall \ell\in\mathcal{L},\qquad
\tfrac{1}{|\mathcal{L}|}\textstyle\sum_{\ell}\Delta_{r,\ell}\geq\tau_{\mathrm{imp}},
\label{eq:promotion-main}
\end{equation}
and it does not increase \texttt{wrong\_format} terminal failures. Resource
statistics and a nominal cost band are recorded, but the cost band was not an
enforced promotion constraint in the reported OOLONG run. The per-stratum
margins $\tau_\ell$ are calibrated to the
harness's own evaluation noise; any one stratum can veto a candidate whose
macro mean improves. When early stopping is enabled, $C_r^\star$'s strata are
evaluated in order, and a completed stratum that already violates its margin
stops validation early.
A promoted single-surface edit becomes $H_r$; a rejection sets
$H_r=H_{r-1}$ and retires
the remaining screened candidates, and the next round mines fresh evidence
against the unchanged incumbent. Validation tasks are never mined, but their
aggregate outcomes inform later proposal history.

\paragraph{Refinement.}
A rejected $C_r^\star$ may still help on the runs where it produced a measurable
answer while failing more often in some other way (for example, malformed
output). If its quality conditional on a measurable answer exceeds the
incumbent's by a configured margin, one bounded repair generation runs (when
enabled; OOLONG-Pairs only in the current study). $C_r^\star$ is mined
on the held-in tasks, the proposer is restricted to patterns with the
configured failure causes, and the single resulting descendant is validated
against the original incumbent under Equation~\ref{eq:promotion-main}. The
branch never recurses and cannot promote without passing the ordinary gate.

\paragraph{Stopping.}
Optimization stops at the round ceiling $T$ or after $p$ consecutive rounds
without promotion, then freezes the final incumbent for separate evaluation.
Appendix~\ref{app:repro} gives the artifact and measurement definitions.

\begin{algorithm*}[t]
\caption{Executed SH-RLM optimization and promotion procedure}
\label{alg:self-harness}
\small
\hrule\vspace{2pt}
\textbf{Require:} fixed model $\mathcal{M}$, initial harness $H_0$, held-in
split $D_{\mathrm{in}}$, held-out split $V$, screening subset $V_s$, evaluator
$\mathcal{E}$, proposal width $k$, round limit $T$, patience $p$.\\[-1pt]
\textbf{Ensure:} frozen final incumbent $H_{T'}$, where $T'\leq T$.\\[-2pt]
\begin{tabular}{@{}r@{\hspace{5pt}}p{.69\textwidth}@{\hspace{8pt}}p{.21\textwidth}@{}}
1 & $z\leftarrow0$ & \emph{consecutive rejections}\\
2 & \textbf{for} $r=1,\ldots,T$ \textbf{do} & \\
3 & \quad $R_r\leftarrow\textsc{Evaluate}(\mathcal{M},H_{r-1},D_{\mathrm{in}},\mathcal{E})$ & held-in executions only\\
4 & \quad $B_r\leftarrow\textsc{MineFailures}(R_r)$ & verifier-grounded evidence\\
5 & \quad $\mathcal{C}_r\leftarrow\textsc{Propose}(\mathcal{M},H_{r-1},B_r,k)$ & independent one-surface edits\\
6 & \quad $\mathcal{C}_r\leftarrow\textsc{Preflight}(\mathcal{C}_r)$ & reject malformed edits\\
7 & \quad \textbf{if} $|\mathcal{C}_r|>1$ \textbf{then} & \\
8 & \qquad $C_r^\star\leftarrow\textsc{TopRanked}(\textsc{Screen}(H_{r-1},\mathcal{C}_r,V_s,\mathcal{E}))$ & screening orders; never promotes\\
9 & \qquad \textbf{if} no candidate clears the screening gate \textbf{then} re-attribute $R_r$ once, propose a fresh pool, and repeat line 8 & bounded fallback\\
10 & \quad \textbf{else} $C_r^\star\leftarrow$ the sole preflighted candidate & skip tournament\\
11 & \quad $(P(H_{r-1}),P(C_r^\star))\leftarrow\textsc{FullValidation}(H_{r-1},C_r^\star,V,\mathcal{E})$ & $v$ attempts per task\\
12 & \quad $\Delta_{r,\ell}\leftarrow\bar F_\ell(C_r^\star)-\bar F_\ell(H_{r-1})$ for every $\ell\in\mathcal{L}$ & all-attempt means\\
13 & \quad $\tau_{\mathrm{imp}}\leftarrow0$ if $r\leq5$, else $.01$ & adaptive guard; round 6 recomputed\\
14 & \quad $a_r\leftarrow[\Delta_{r,\ell}\geq-\tau_\ell\ \forall\ell]\wedge[|\mathcal{L}|^{-1}\sum_\ell\Delta_{r,\ell}\geq\tau_{\mathrm{imp}}]$ & quality gate\\
15 & \quad $a_r\leftarrow a_r\wedge[\mathrm{WF}(C_r^\star)\leq\mathrm{WF}(H_{r-1})]$ & wrong-format gate\\
16 & \quad \textbf{if} $a_r$ \textbf{then} $H_r\leftarrow C_r^\star$; $z\leftarrow0$; record \textsc{Promote} & incumbent changes\\
17 & \quad \textbf{else if} the refinement trigger fires \textbf{then} $C_r^\dagger\leftarrow\textsc{Refine}(C_r^\star,B_r)$ once; repeat lines 11--15 for $C_r^\dagger$ and set $(H_r,z)$ to $(C_r^\dagger,0)$ if it passes, otherwise $(H_{r-1},z+1)$ & no recursive refinement\\
18 & \quad \textbf{else} $H_r\leftarrow H_{r-1}$; $z\leftarrow z+1$; record \textsc{Reject} & incumbent unchanged\\
19 & \quad Persist candidates, attempts, summaries, decisions, hashes, and any configuration intervention & complete audit trail\\
20 & \quad \textbf{if} $z\geq p$ \textbf{then break} & patience stopping\\
21 & \textbf{end for} & \\
22 & \textbf{return} the frozen final incumbent $H_{T'}$ & separate final evaluation\\
\end{tabular}
\vspace{2pt}\hrule
\end{algorithm*}
